% ===========================================================================
% 13. Discussion
% ===========================================================================
\section{Discussion}
\label{sec:discussion}

\subsection{Mitigations Work for Framing, Not for Injection}
\label{sec:discuss-asymmetry}

The most striking pattern in our data is the \emph{asymmetry} of
mitigation effectiveness across attack types. Semantic manipulation
(authority framing: 100\%$\to$0\%), human-in-the-loop attacks (all
three scenarios: 40--60\%$\to$0\%), and cross-contamination
(100\%$\to$0\%) are fully mitigated by the hardened configuration.
These are attacks where the adversarial signal operates at the
\emph{semantic} level---framing, selective reporting, social
engineering---and can be countered by prompt-level defenses that
instruct the agent to be skeptical.

In contrast, content injection via CSS invisible text and dynamic
cloaking (100\% in both conditions) completely resists all
mitigations. These attacks exploit a \emph{structural} gap: LLMs
process raw text without awareness of rendering context. No amount of
prompt engineering can fix this; architectural defenses (e.g., content
rendering pipelines that strip non-visible elements before LLM
processing) are required.

\subsection{Regression Cases: Iatrogenic Mitigations}
\label{sec:discuss-regression}

Two scenarios exhibit \emph{regression}---the hardened condition
performs worse than baseline:

\begin{enumerate}[nosep,leftmargin=*]
  \item \textbf{Misleading forms (BC-2)}: 0\%$\to$100\%
    ($d = \infty$, $p = 0.033$). The model naturally resists this
    attack, but the mitigation suite introduces the vulnerability.
    The security-hardened system prompt may alter form-processing
    behavior in ways that make the agent \emph{more} susceptible.

  \item \textbf{Ranking manipulation (CS-2)}: 80\%$\to$100\%
    ($d = -0.63$). The \mitigation{rag-integrity} module's
    cross-reference validation may inadvertently amplify manipulated
    rankings by treating them as authoritative signals.
\end{enumerate}

These regressions constitute an \emph{iatrogenic} phenomenon:
defenses that cause harm. This is a critical finding for the security
community---mitigation suites must be evaluated holistically, not
just on their target scenarios, as they can create new attack surfaces.

\subsection{Natural Resistance of GPT-4o-mini}
\label{sec:discuss-natural}

Of 22 scenarios, \textbf{12 show 0\% baseline trap success}
(SM-2--SM-4, BC-3--BC-4, SY-1--SY-3, and three partially at 40\%),
indicating that GPT-4o-mini's RLHF training already provides
substantial protection against many attack vectors. The systemic
category is particularly notable: all three scenarios achieve 0\%
in both conditions, suggesting that protocol-level cascade attacks
may require more sophisticated model targeting or larger agent
topologies.

\subsection{Statistical Limitations}
\label{sec:discuss-stats}

Our $n=5$ design has important limitations:

\begin{itemize}[nosep,leftmargin=*]
  \item \textbf{No individual comparison reaches Bonferroni
    significance} ($\alpha_{\text{corrected}} = 0.00227$). The lowest
    $p$-values are 0.033 (SM-1, CS-4, BC-2), well above the
    corrected threshold. We interpret results through effect sizes
    rather than NHST.

  \item \textbf{The hardened condition confounds} mitigation modules
    with the security-hardened prompt. A 2$\times$2 factorial
    (modules $\pm$ prompt $\pm$) would disambiguate these effects.

  \item \textbf{Single model.} Our results characterize GPT-4o-mini
    only. Frontier models may show different vulnerability profiles;
    multi-model comparison is a priority for follow-up work.
\end{itemize}

\subsection{The 20pp Gap: Modest but Meaningful}
\label{sec:discuss-gap}

The overall 20 percentage-point reduction (45.5\%$\to$25.5\%) is
modest. A quarter of all trap attempts still succeed under full
hardening. This suggests that prompt-level and module-level defenses
are necessary but insufficient; \emph{architectural} defenses---content
rendering pipelines, embedding-level anomaly detection, protocol-level
authentication---are needed to close the remaining gap.

\subsection{Protocol-Level Implications}
\label{sec:discuss-protocol}

Although systemic attacks showed 0\% success in our study, the A2A
protocol's lack of message authentication, content integrity, and
trust frameworks remains a structural concern. We recommend that
future protocol versions incorporate mandatory message signing
(e.g., JWS~\cite{rfc7515}), content integrity hashes, and trust
scoring.

\subsection{Limitations}
\label{sec:discuss-limitations}

\begin{enumerate}[nosep,leftmargin=*]
  \item \textbf{Synthetic environments.} Our trap scenarios target
    controlled measurement, not realism.
  \item \textbf{Static adversaries.} Our adversaries do not adapt to
    the model or defense configuration.
  \item \textbf{No ablation or compound data.} This run covers only
    baseline vs.\ hardened; per-module and compound analyses are
    planned.
  \item \textbf{Model versioning.} Results reflect a specific
    GPT-4o-mini snapshot and may not generalize to future versions.
  \item \textbf{HitL validation.} HL results measure biased
    \emph{output}, not human \emph{deception}.
  \item \textbf{Tool-use attacks not covered.} Following the DeepMind
    taxonomy scope, we evaluate environmental/contextual attacks and
    do not cover tool-use/function-calling attack surfaces
    (cf.\ AgentDojo~\cite{debenedetti2024agentdojo}).
\end{enumerate}

\subsection{Broader Impact}
\label{sec:discuss-impact}

This work has dual-use implications. By publishing adversarial
scenarios, we enable defenders to test and harden their systems---but
we also provide a blueprint for potential attackers. We believe the
defensive value outweighs the risk, as the attack techniques we
operationalize are already described in the public DeepMind
taxonomy~\cite{yuntao2025agenttraps}. Our contribution is the empirical
measurement and defense evaluation, not the attack invention.
